
\begin{table*}[!t]
\begin{tabular}{lp{0.82\textwidth}}
\toprule
Question & 
Compound X, which has the following IR and NMR data, reacts with red phosphorus and HI. Identify the final product. IR: 3400–2500 cm-1, 1720 cm-1, 1610 cm-1, 1450 cm-1 1H NMR: 10.5 ppm (bs, 1H), 8.0 ppm (d, 2H), 7.2 ppm (d, 2H), 2.9 ppm (m, 1H), 1.7 ppm (m, 2H), 1.4 ppm (d, 3H), 0.9 ppm (t, 3H).  A) 4-(sec-butyl)benzoic acid B) 1-isobutyl-4-methylbenzene C) 2-(4-ethylphenyl)propanoic acid D) 1-(sec-butyl)-4-methylbenzene 
 \\
\midrule 
Answer & D \\
\midrule
Instruct & Let's solve this step by step:  1) First, let's analyze the IR data

...

5) Looking at the options: - Only option A has the carboxylic acid group - The other options don't have carboxylic acid groups - The NMR pattern matches a para-substituted benzene with a sec-butyl group  The original compound is 4-(sec-butyl)benzoic acid, which upon reduction with red phosphorus and HI gives 1-(sec-butyl)-4-methylbenzene.  \textcolor{blue}{ANSWER: D}
\\
\midrule
Thinking &  Okay, let's try to figure out this problem step by step. So, we have Compound X with some IR and NMR data, and it reacts with red phosphorus and HI. 

...

Conclusion: The correct final product after the reaction of Compound X with red phosphorus and HI is 1-(sec-butyl)-4-methylbenzene.  \textcolor{blue}{ANSWER: D}
 \\
\midrule
MI-0.2 & Let's analyze the data step by step:

...

After reduction, it becomes 1-(isobutyl)-4-methylbenzene.  This matches option B: 1-isobutyl-4-methylbenzene.  \textcolor{red}{ANSWER: B} 
\\
\midrule
MI-0.8 &  First, I need to analyze the IR and NMR data to figure out what Compound X is.

...

  The final product is identified as D.  \textcolor{blue}{ANSWER: D}
\\
\bottomrule
\end{tabular}
\caption{
One case from the GPQA-Diamond benchmark.
We skip the details.
MI-0.2 fails while all others can solve it.
Also, the style for the beginning of MI-0.2 is similar to the Instruct model.
}
\label{appdx:case2}
\end{table*}